Un protó en repòs és accelerat en el sentit positiu de l'eix $x$ fins a assolir una velocitat d'$1,00\times 10^{5}\ \mathrm{m\,s^{-1}}$. Aleshores, penetra en un espectròmetre de masses on hi ha un camp magnètic $\vec{B} = 1,00\times 10^{-2}\ \mathrm{T}\,\vec{k}$.

\begin{apartat}{1}
Calculeu la força (mòdul, direcció i sentit) que actua sobre el protó.
\end{apartat}

\begin{apartat}{1}
Calculeu el camp magnètic (mòdul, direcció i sentit) tal que, si entra un electró amb la mateixa velocitat en l'espectròmetre, segueixi la mateixa trajectòria que el protó.
\end{apartat}

\dades{%
  Càrrega elemental, $e = 1,60\times 10^{-19}\ \mathrm{C}$.\\
  Massa del protó, $m_\mathrm{p} = 1,67\times 10^{-27}\ \mathrm{kg}$.\\
  Massa de l'electró, $m_\mathrm{e} = 9,11\times 10^{-31}\ \mathrm{kg}$.\\
  $k = \dfrac{1}{4\pi\varepsilon_0} = 8,99\times 10^{9}\ \mathrm{N\,m^2\,C^{-2}}$.}
